% answers_short.tex —— 2004 年·普通高等学校招生考试·广东卷 · 简短答案速查
% 由 tools/make_short_answers.py 自 answers.tex 的 \gksolhead/\gkanswer 提取，
% 并与 2004/2004_数学真题_广东_答案要点.md 逐题比对后生成，勿手改（改源 MD 后重跑）。

\begin{gknotice}
\textbf{使用说明}：本节只给最终答案便于快速核对，完整推导见解析版。答案为编写者独立推导，\textbf{非官方标准答案}。
\end{gknotice}

\gkpart{选择题}
\begin{center}\begin{tabular}{ccccccccccccc}
\toprule
题号 & 1 & 2 & 3 & 4 & 5 & 6 & 7 & 8 & 9 & 10 & 11 & 12 \\
\midrule
答案 & C & A & C & A & B & D & D & A & A & B & D & B \\
\bottomrule
\end{tabular}\end{center}

\gkpart{填空题}
\begin{itemize}[leftmargin=2.6em]
  \item[\textbf{13.}] $\dfrac{5}{7}$
  \item[\textbf{14.}] $-2\mathrm{i}$
  \item[\textbf{15.}] $\dfrac{PA' \cdot PB' \cdot PC'}{PA \cdot PB \cdot PC}$
  \item[\textbf{16.}] $(\mathrm{e}^x + 1)^2 - 1$（或 $\mathrm{e}^{2x} + 2\mathrm{e}^x$），$x \in \mathbb{R}$
\end{itemize}

\gkpart{解答题}
\begin{itemize}[leftmargin=2.6em]
  \item[\textbf{17.}] $\alpha = \dfrac{2\pi}{3}, \beta = \dfrac{4\pi}{3}, \gamma = \dfrac{8\pi}{3}$ 或 $\alpha = \dfrac{4\pi}{3}, \beta = \dfrac{8\pi}{3}, \gamma = \dfrac{16\pi}{3}$
  \item[\textbf{18.}] （1）二面角 $C - DE - C_1$ 的正切值为 $\dfrac{\sqrt{2}}{2}$；（2）直线 $EC_1$ 与 $FD_1$ 所成角的余弦值为 $\dfrac{\sqrt{21}}{14}$
  \item[\textbf{19.}] （1）证明见解析；（2）切线与两坐标轴正向所围成的三角形面积表达式为 $S = \dfrac{1}{2}(2 - x_0)^2$
  \item[\textbf{20.}] 巨响发生在中心的西偏北 $45^\circ$、距中心 $680\sqrt{10}\text{ m}$ 处（以中心为原点，正东为 $x$ 轴正方向，正北为 $y$ 轴正方向时，坐标为 $(-680\sqrt{5}, 680\sqrt{5})$）
  \item[\textbf{21.}] （1）当整数 $m \leqslant 1$ 时，$f(x) \geqslant 0$；（2）证明见解析
  \item[\textbf{22.}] 直线 $l$ 的方程为 $x = \pm \dfrac{25\sqrt{241}}{241}$，或 $y = \pm \dfrac{16}{13}$，或 $y = \pm \dfrac{16}{25}x$（即 $16x \pm 25y = 0$）
\end{itemize}
